% Generated draft: requires mathematical and source review.
\documentclass[11pt]{article}
\usepackage[margin=25mm]{geometry}
\usepackage{fontspec}
\setmainfont{Latin Modern Roman}
\usepackage{amsmath,amssymb,mathtools,unicode-math}
\setmathfont{Latin Modern Math}
\usepackage{enumitem,needspace,xurl,hyperref}
\setlength{\parindent}{0pt}
\setlength{\parskip}{5pt}
\setlength{\emergencystretch}{3em}
\newcommand{\N}{\mathbb{N}}
\newcommand{\Z}{\mathbb{Z}}
\newcommand{\Q}{\mathbb{Q}}
\newcommand{\R}{\mathbb{R}}
\newcommand{\C}{\mathbb{C}}
\newcommand{\F}{\mathbb{F}}
\newcommand{\A}{\mathbb{A}}
\newcommand{\PP}{\mathbb{P}}
\newcommand{\E}{\mathbb{E}}
\newcommand{\eps}{\varepsilon}
\newcommand{\dd}{\,\mathrm{d}}
\begin{document}
\section*{UnsolvedMath ID: 30000469}
\textbf{Source problem code:} catalogue entry
\subsection{Definitions and mathematical statement}
The source asks whether large-cardinal axioms are compatible with weaker forms of Condensation. In the fine-structural context, a condensation principle typically relates the transitive collapse of an appropriate elementary substructure of a hierarchy level to a level or initial segment of the same hierarchy. Its content depends on the hierarchy, predicates, substructure size, and collection of substructures for which the conclusion is required. None of these data is specified in the record.

For a fixed large-cardinal assertion \(\mathsf L\) and a fixed weak condensation principle \(\mathsf C\), the intended compatibility question would concern
\[
\operatorname{Con}(\mathrm{ZFC}+\mathsf L+\mathsf C),
\]
and the hypotheses from which that consistency follows. A same-strength positive relative-consistency result might have the form
\[
\operatorname{Con}(\mathrm{ZFC}+\mathsf L)
\ \Longrightarrow\ 
\operatorname{Con}(\mathrm{ZFC}+\mathsf L+\mathsf C),
\]
but the extract does not require or establish that particular bound; stronger starting assumptions or different preserved large cardinals may be involved. Compatibility is not the assertion that a large-cardinal axiom implies condensation in every universe.

The source does not identify \(\mathsf L\) or \(\mathsf C\), or say whether it asks for all large-cardinal levels, one example, or optimal strength for a particular weak principle. The supplied literature assessment checked 2026-08-21 distinguishes stationary and local-club condensation forms, reports consistency results with large cardinals for such forms, and says that the full consistency-strength question for the unspecified weaker forms remains unresolved. These named variants are not interchangeable, and their exact definitions and compatible cardinal hypotheses are not reproduced. The item consequently remains a family of relative-consistency questions, not a single fully specified theorem.
\subsection{Short English statement}
Determine which weak condensation principles can hold together with which large-cardinal axioms, and at what consistency strength. The source names neither the precise weak principle nor the cardinal strength, so the reported partial results cannot be promoted to a universal compatibility theorem.
\subsection{Sources}
\begin{itemize}
\item Local UnsolvedMath registry, permanent ID 30000469, catalogue entry, \emph{Weak Condensation with Large Cardinals}.
\item \emph{Feinstrukturtheorie und Innere Modelle} (2007), Oberwolfach Reports, pp. 1215--1244, DOI \url{https://doi.org/10.4171/owr/2006/20}.
\item Friedman--Holy, supplied title \emph{Condensation and large cardinals}: \url{https://eudml.org/doc/286512}.
\item Supplied literature triage checked 2026-08-21 reports stationary and local-club partial consistency results without their exact definitions or hypotheses.
\end{itemize}
\end{document}
